\documentclass{beamer}

\usepackage{cuhkbeamer}

\begin{document}

% ====================== Slide 1: TITLE SLIDE (REVISED) ======================
\begin{frame}
  \title{Synthesis \& Capstone Group Presentations}
  \author{Chiu Yu KO}
  \institute{Low Altitude Economics}
  \date{Week 11}
  \titlepage
\end{frame}

% ====================== Slide 2: FORMAT & OBJECTIVES ======================
\begin{frame}[t]
\frametitle{Learning Objectives}
By the end of this class, you should be able to:
\begin{itemize}
\item \textbf{Integrate} demand, capacity, unit economics, infrastructure, contracts, regulation, and sustainability into one coherent strategy.
\item \textbf{Check} that assumptions, units, and causal links remain consistent across the complete business model.
\item \textbf{Evaluate} viability using cash flow, NPV, breakeven, downside risk, and non-financial decision criteria.
\item \textbf{Present} a concise 12-minute recommendation supported by transparent evidence and sensitivity analysis.
\item \textbf{Defend} the proposal by explaining limitations, decision gates, and conditions under which the recommendation should change.
\end{itemize}
\end{frame}

% ====================== Slide 3: ICEBREAKER ======================
\begin{frame}[t]
\frametitle{Seeing the Whole System}
\textbf{A strong capstone does more than collect ten weeks of separate calculations.}
\begin{itemize}
\item Demand determines the activity the system may serve.
\item Infrastructure and regulation constrain reliable capacity.
\item Capacity, utilization, and contracts determine revenue and unit cost.
\item Environmental and social impacts affect operating conditions, mitigation cost, and legitimacy.
\end{itemize}
\vspace{0.3cm}
\textbf{Synthesis means tracing how one changed assumption affects the rest of the model.}
\begin{itemize}
\item \textbf{Quick poll (2 min):} Which assumption in your project creates the largest cascade if it changes?
\end{itemize}
\end{frame}


% ====================== Session 1 Title Card ======================
\begin{frame}[c]
  \begin{center}
    \huge\color{cuhkpurple}\textbf{Session 1}\\
    \vspace{0.5cm}
    \Large Integrated Strategy Synthesis \& The Master Viability Equation
  \end{center}
\end{frame}

% ====================== Slide 4: CONCEPT PRIMER 1 ======================
\begin{frame}[t]
\frametitle{Concept Primer: The Integration Challenge}
Business-model components are connected through quantities, prices, capacity, cost, timing, risk, and stakeholder responses.
\begin{itemize}
\item A restriction on operating hours reduces reliable capacity.
\item Lower capacity may reduce served demand and revenue while increasing fixed cost per served unit.
\item A price response may further change customer choice and load factor.
\item A mitigation measure may increase cost but preserve operating access or reduce disruption.
\end{itemize}
\textbf{Integration rule:} Every important assumption should appear once in the correct model channel, with no contradictory units or double counting.
\end{frame}

% ====================== Slide 5: STANDALONE FORMULA ======================
\begin{frame}[t]
\frametitle{Applied Tool: The Viability Test}

% Tighten vertical math spacing inside blocks
\setlength{\abovedisplayskip}{3pt}
\setlength{\belowdisplayskip}{3pt}

\begin{block}{Step 1: How Much Can We Actually Serve?}
\[
\text{Served Activity} = \min\{ \text{Customer Demand}, \text{Reliable Capacity} \}
\]
\end{block}

\vspace{-0.1cm}

\begin{block}{Step 2: Does the Service Cover Its Costs?}
\[
\text{Operating Profit} = (\text{Price} - \text{Variable Cost}) \times \text{Served Activity} - \text{Fixed Cost}
\]
\end{block}

\vspace{0.15cm}

\begin{itemize}
  \setlength{\itemsep}{2pt}
  \item Demand may be limited by price, travel time, or customer acceptance.
  \item Capacity may be limited by aircraft, infrastructure, weather, or regulation.
  \item The business is operationally viable only if served activity generates enough contribution to cover fixed cost.
\end{itemize}

\end{frame}


% ====================== Slide 6: BACK-OF-THE-ENVELOPE ======================
\begin{frame}[t]
\frametitle{Back-of-the-Envelope: The Effect of a Capacity Shock}

% Control math whitespace inside blocks
\setlength{\abovedisplayskip}{2pt}
\setlength{\belowdisplayskip}{2pt}

\small
\textbf{Assumptions:} Customer demand is 100 flights/day at HK\$600 per flight. Variable cost is HK\$100 per flight, and fixed cost is HK\$40,000/day.

\vspace{0.15cm}

\begin{columns}[T]
  \begin{column}{0.48\textwidth}
    \begin{block}{Base Case}
      \textbf{Step 1: Served Activity}
      \[
      \min\{100, 100\} = 100
      \]
      \vspace{-0.1cm}
      \textbf{Step 2: Operating Profit}
      \[
      \begin{aligned}
        & (600 - 100) \times 100 - 40{,}000 \\
        ={} & \mathbf{HK\$10{,}000/\text{day}}
      \end{aligned}
      \]
    \end{block}
  \end{column}

  \begin{column}{0.48\textwidth}
    \begin{block}{Capacity Shock}
      \textbf{Step 1: Served Activity}
      \[
      \min\{100, 80\} = 80
      \]
      \vspace{-0.1cm}
      \textbf{Step 2: Operating Profit}
      \[
      \begin{aligned}
        & (600 - 100) \times 80 - 40{,}000 \\
        ={} & \mathbf{HK\$0/\text{day}}
      \end{aligned}
      \]
    \end{block}
  \end{column}
\end{columns}

\vspace{0.2cm}

\footnotesize
\textit{Takeaway: The capacity shock reduces served activity even though customer demand is unchanged. Because fixed cost remains, the service falls from a daily profit of HK\$10,000 to breakeven.}

\vfill
{\tiny Illustrative assumptions; no claim about annual Hong Kong weather downtime.}
\end{frame}

% ====================== Slide 7: CONCEPT PRIMER 2 ======================
\begin{frame}[t]
\frametitle{Concept Primer: Integrated Strategy Template}
A coherent capstone should connect six decision blocks:
\begin{enumerate}
\item \textbf{Service and stakeholder problem}: What need is addressed, for whom, and compared with what alternative?
\item \textbf{Demand and capacity}: How many units are demanded and reliably serviceable?
\item \textbf{Economics and finance}: What are contribution, breakeven activity, cash flow, NPV, and downside exposure?
\item \textbf{Infrastructure and operations}: Which resources are binding, and how is reliability maintained?
\item \textbf{Partnership and regulation}: Who performs each role, bears each risk, and grants each permission?
\item \textbf{Sustainability and distribution}: What impacts arise, who bears them, and how are they measured and mitigated?
\end{enumerate}
\end{frame}

% ====================== NEW SLIDE 7.5: INTEGRATED SYNTHESIS EXAMPLE ======================
\begin{frame}[t,shrink=10]
\frametitle{Integrated Synthesis Example: Airport-Corridor Screening}
\textbf{Illustrative chain of reasoning:}
\begin{itemize}
\item \textbf{Problem}: Test whether a time-sensitive airport segment values the proposed door-to-door service.
\item \textbf{Demand}: Estimate customer choice at alternative fares using generalized cost and capacity constraints.
\item \textbf{Operations}: Define aircraft, seats, load factor, operating days, turnaround, and serviceable capacity.
\item \textbf{Economics}: Calculate contribution per passenger-trip, fixed cost, breakeven activity, NPV, and downside cash need.
\item \textbf{Regulation}: Identify applicable permissions and use evidence-based Sandbox X or AC-013 milestones only where relevant.
\item \textbf{Impact}: Compare life-cycle emissions, noise exposure, privacy controls, access, and public-service effects.
\end{itemize}
\textbf{Decision:} Recommend proceed, stage, redesign, relocate, or reject, with explicit conditions that would change the recommendation.
\end{frame}




% ====================== NEW SLIDE 7.6: RISK REGISTER & ROADMAP ======================
\begin{frame}[t]
\frametitle{Applied Tool: Risk Register and Decision-Gated Roadmap}

\begin{center}\footnotesize
\setlength{\tabcolsep}{4pt}
\renewcommand{\arraystretch}{1.15}
\begin{tabularx}{\linewidth}{>{\raggedright\arraybackslash}p{1.8cm}|>{\raggedright\arraybackslash}X|>{\raggedright\arraybackslash}X|>{\raggedright\arraybackslash}X}
\textbf{Risk} & \textbf{Model channel} & \textbf{Response} & \textbf{Decision trigger}\\
\hline
Demand & Choice share and price & Pilot and staged capacity & Minimum verified demand\\
Weather & Serviceable capacity & Procedures and resilience & Reliability threshold\\
Regulation & Scope and timing & Evidence plan and review & Permission or milestone\\
Technology & Cost and availability & Test, redundancy, supplier plan & Performance threshold\\
Community impact & Route and operating conditions & Mitigation and engagement & Exposure and complaint KPI\\
\end{tabularx}
\end{center}

\vspace{0.15cm}
\footnotesize
\textit{Do not assume fixed Sandbox X approval dates. Under AC-013, BVLOS permission within the regulatory-sandbox context is operation-specific and requires supporting documents, flight and test plans, risk assessment, and mitigations.}
\end{frame}

% ====================== Slide 8: SESSION 2 TITLE ======================
\begin{frame}[c]
  \begin{center}
    \huge\color{cuhkpurple}\textbf{Session 2}\\
    \vspace{0.5cm}
    \Large Presentation Skills, Practitioner Frameworks \& Rehearsal
  \end{center}
\end{frame}

% ====================== Slide 9: CONCEPT PRIMER 3 ======================
\begin{frame}[t]
\frametitle{Concept Primer: The Decision-Maker's Lens}
Assessors may represent investors, operators, regulators, public agencies, customers, or affected communities. Their objectives are not identical.
\begin{itemize}
\item \textbf{Operational credibility}: Can the service be delivered safely and reliably?
\item \textbf{Commercial viability}: Are demand, contribution, cash flow, and financing assumptions credible?
\item \textbf{Regulatory feasibility}: Is the proposed pathway specific to the aircraft and operation?
\item \textbf{Public value}: Are benefits, costs, access, and distribution addressed?
\item \textbf{Decision quality}: Does the recommendation acknowledge uncertainty, alternatives, and conditions for revision?
\end{itemize}
\textbf{A strong presentation answers all five without pretending that uncertainty has disappeared.}
\end{frame}

% ====================== Slide 10: STANDALONE FORMULA ======================
\begin{frame}[t]
\frametitle{Applied Tool: A 12-Minute Decision Pitch}
\begin{itemize}
\item \textbf{Decision and value proposition, 2 minutes}: State the problem, recommendation, customer, use case, and principal evidence.
\item \textbf{Integrated quantitative case, 4 minutes}: Show demand, reliable capacity, contribution, breakeven, NPV, and one clear sensitivity.
\item \textbf{Implementation and governance, 3 minutes}: Explain infrastructure, partners, regulatory pathway, responsibilities, and decision gates.
\item \textbf{Risk and impact, 2 minutes}: Present the major downside scenario, mitigation, residual risk, sustainability, and distributional effects.
\item \textbf{Closing decision, 1 minute}: State the requested action, next milestone, funding or approval need, and conditions for stopping.
\end{itemize}
\textit{Use formulas only where they clarify the decision. The objective is a defensible recommendation, not the largest number of calculations.}
\end{frame}

% ====================== Slide 11: BACK-OF-ENVELOPE (WIN vs LOSE) ======================
\begin{frame}[t,shrink=13]
\frametitle{Back-of-the-Envelope: Weak and Strong Evidence}
\begin{columns}[T]
\begin{column}{0.48\textwidth}
\begin{block}{Weak Case}
\begin{itemize}
\item Uses total market size as forecast demand
\item Confuses passenger-trips, flights, and seats
\item Applies repeated percentage haircuts
\item Reports expected profit without downside risk
\item Assumes approval, acceptance, and insurance outcomes
\end{itemize}
\end{block}
\end{column}
\begin{column}{0.48\textwidth}
\begin{block}{Strong Case}
\begin{itemize}
\item Defines the market and competing alternative
\item Checks demand against reliable capacity
\item Maps each constraint to one model channel
\item Reports NPV, probability of loss, and cash need
\item States evidence, limitations, and decision gates
\end{itemize}
\end{block}
\end{column}
\end{columns}
\textit{A strong analysis can still recommend rejection or delay. Credibility is more important than producing a favourable answer.}
\end{frame}

% ====================== Slide 12: SESSION 3 TITLE ======================
\begin{frame}[c]
  \begin{center}
    \huge\color{cuhkpurple}\textbf{Session 3}\\
    \vspace{0.5cm}
    \Large Capstone Group Presentations + Practitioner Feedback
  \end{center}
\end{frame}

% ====================== Slide 13: PRESENTATION LOGISTICS ======================
\begin{frame}[t]
\frametitle{Presentation Schedule and Ground Rules}
\begin{itemize}
\item \textbf{Timing}: 12 minutes to present and 3 minutes for questions. Rehearse to finish without rushing the recommendation.
\item \textbf{Team contribution}: Every member should speak and understand the central assumptions and results.
\item \textbf{Traceability}: Cite important data and distinguish sourced evidence from classroom assumptions.
\item \textbf{Visual discipline}: Use readable charts, tables, sensitivity results, and a decision roadmap rather than dense text.
\item \textbf{Unit consistency}: Distinguish passengers, passenger-trips, flights, movements, operating days, and annual values.
\item \textbf{Professional judgment}: Explain uncertainty and limitations directly. Do not present assumptions as current Hong Kong facts.
\end{itemize}
\end{frame}

% ====================== Slide 14: FEEDBACK GUIDELINES ======================
\begin{frame}[t]
\frametitle{Feedback Guidelines for the Audience}
When evaluating another team, provide one strength, one material concern, and one actionable improvement.
\begin{itemize}
\item \textbf{Problem and recommendation}: Is the decision clear and tied to a real stakeholder need?
\item \textbf{Model consistency}: Do demand, capacity, price, cost, and timing use compatible units?
\item \textbf{Evidence}: Are important facts sourced and uncertain inputs varied?
\item \textbf{Implementation}: Are partners, permissions, responsibilities, and milestones plausible?
\item \textbf{Risk}: Does the analysis identify residual risk, probability of loss, and liquidity needs?
\item \textbf{Impact}: Are safety, privacy, noise, climate, access, and distribution considered proportionately?
\end{itemize}
\end{frame}

% ====================== Slide 15: EXERCISE TEMPLATE ======================
\begin{frame}[t,shrink=10]
\frametitle{Visual Template: A 12-Slide Executive Deck}
\begin{enumerate}
\item Decision and executive summary
\item Stakeholder problem and current alternative
\item Service design and operating concept
\item Market definition and customer choice
\item Reliable capacity and infrastructure
\item Unit economics and breakeven
\item Cash flow, NPV, and funding requirement
\item Partnership, payment, and risk allocation
\item Regulatory pathway and evidence gates
\item Downside scenario and sensitivity results
\item Sustainability and distributional impacts
\item Recommendation, milestones, and stop or expand conditions
\end{enumerate}
\textit{The deck order may change, but the logic should move from the decision to evidence, implementation, risk, and action.}
\end{frame}

% ====================== Slide 16: KEY TAKEAWAYS ======================
\begin{frame}[t]
\frametitle{Key Takeaways}
\begin{itemize}
\item Synthesis requires a connected model, not a collection of independent weekly calculations.
\item Served activity is limited by both customer demand and reliable capacity.
\item Project viability should be assessed through contribution, cash flow, NPV, downside risk, and non-financial criteria.
\item Constraints must enter through explicit model channels without double counting.
\item A credible recommendation states evidence, limitations, residual risks, and conditions for changing course.
\item Clear presentation supports judgment by making assumptions and trade-offs visible rather than hiding uncertainty.
\end{itemize}
\end{frame}

% ====================== Slide 17: ASSIGNMENT ======================
\begin{frame}[t]
\frametitle{Week 11 Deliverables}
\textbf{Capstone Presentation and Executive Memo (30\% of course grade)}
\begin{itemize}
\item \textbf{Due}: End of Week 11
\item \textbf{Presentation}: 12 minutes plus 3 minutes for questions
\item \textbf{Memo}: Maximum five pages excluding references and clearly labelled appendices
\item \textbf{Required content}:
\begin{enumerate}
\item integrated recommendation covering service, demand, capacity, economics, infrastructure, partnership, regulation, and impacts;
\item consistent quantitative model including breakeven and NPV or an explained alternative;
\item downside analysis reporting probability of loss or a comparable risk measure and liquidity need;
\item risk register with owners, responses, residual risks, and decision triggers; and
\item phased implementation roadmap linked to evidence rather than assumed approval dates.
\end{enumerate}
\item \textbf{Grading focus}: Integration and model consistency (35\%), evidence and feasibility (25\%), risk and implementation (20\%), impact and governance (10\%), presentation clarity (10\%).
\end{itemize}
\end{frame}

% ====================== Slide 18: PREVIEW WEEK 12 ======================
\begin{frame}[t]
  \frametitle{Preview of Week 12}
  
  \begin{center}
      \Large \textbf{Review, Quantitative Applications \& Final Exam}
  \end{center}
  
  \vspace{1cm}
  \textbf{Preview question for next week:} \\
  \textit{How will you demonstrate this level of integrated economic judgment across all 11 weeks under the time pressure of a final exam?}
\end{frame}

% ====================== Slide 19: THANK YOU ======================
\begin{frame}[t]
  \frametitle{Thank You + Q\&A}
  
  \textbf{Pre-Class Readings Reminder (for Week 12):}
  \begin{itemize}
      \item \textit{Review all past reading bundles.}
      \item \textit{Prepare questions for the final exam review session.}
  \end{itemize}
  
  \vspace{1cm}
  \begin{center}
      \textbf{Questions? Final Presentation Feedback?} \\
      \small (Final 15 minutes reserved for extended Q\&A)
  \end{center}
\end{frame}

\end{document}